\section{Results}
\label{sec:results}
\subsection{Hardware Performance}
\label{subsec:hardware_performance}

\subsubsection{Thermal Analysis}
\label{subsec:thermal_analysis}
The Light Tube V3 and Spotlight variants maintain safe operating temperatures.
The aluminum profile of the Tube acts as an effective passive heatsink, while the Spotlight's active cooling is sufficient.
No critical heat buildup occurs at maximum brightness.
The Panel variant overheats.
Direct active cooling of the PCB is insufficient without a dedicated heatsink.
Future designs require a heatsink and fan.
The five addressable sections provide limited utility.
Replacing them with a single section of high-power, MOSFET-controlled LEDs is recommended to eliminate the quiescent current draw of addressable ICs.

\subsubsection{Power Stability}
\label{subsec:power_stability}
The 5V rail remains stable under full logic load (Art-Net and Display active).
LED strip brightness does not affect the 5V rail.
Addressable strips (WS2815) draw power from the 12V/24V input, bypassing the 5V buck converter.
MOSFET gate switching current is negligible.

\subsection{Software Performance}
\label{subsec:software_performance}
\subsubsection{Responsiveness}
\label{subsec:responsiveness}
Art-Net over Wi-Fi exhibits visible latency.
It is unsuitable for live music synchronization.
Wired DMX512 shows no perceptible latency and is validated for live performance.
The DMX driver occasionally fails upon bus connection, requiring a microcontroller reset.
Post-reset operation is stable.

\subsection{Overall System Functionality for Each Product}
\label{subsec:overall_system_functionality}
\subsubsection{Light Tube V2 (PWM Control)}
\label{subsubsec:light_tube_v2_results}
The V2 variant is considered deprecated, mainly due to the enclosure design. 
Integrating the controller into the end cap meant there were no connection points for building "light sculptures".
Also, controlling the whole tube as just one segment via MOSFETs felt visually limited.

\subsubsection{Light Tube V3 (RGB IC Control)}
\label{subsubsec:light_tube_v3_results}
The V3 is a solid improvement. 
The addressable LED strips provide much better visuals.
The "Screw it!" connection design is a highlight, offering the support needed for modular, multi-tube builds.

\subsubsection{Panel Light}
\label{subsubsec:panel_light_results}
The Panel functions, but you have to be careful with high brightness levels. 
The internal logic—sending I2C data to control addressable LEDs—ended up being overly complex for what it achieves.

\subsubsection{Spot Light}
\label{subsubsec:spot_light_results}
The Spotlight works well. The lens focuses the single COB LED effectively, and the 3D-printed diffusion cube blends the colors nicely.
The aluminum heatsink and fan keep it cool enough, and having just one control segment is exactly what you need for a spot.

\subsection{Enclosure Quality and Durability}
\label{subsec:enclosure_quality}
The 3D-printed components have proven to be robust. 
They hold up well under normal use and provide good overall stability for the fixtures. 
The heat-set inserts are a key factor here; they offer enough strength for the "Screw it!" connection system to work reliably without stripping or pulling out. 
Although the parts haven't been purposely tested to failure, no mechanical failures have occurred during assembly or operation so far.